\section{Experimental Methodology}
\label{sec:methodology}

This section describes how \benchmark is evaluated once a dataset release has been frozen. The benchmark definition in \secref{sec:benchmark-design} fixes the meaning of an instance, the syntactic and executable call graphs, and the canonical trace. The methodology below focuses on the experimental protocol: which generated instances are retained, how models are served, how prompts are constructed, how outputs are parsed, and how scores are analyzed. The final dataset parameters are still open and are therefore recorded as controlled placeholders rather than fixed claims.

\subsection{Dataset Release}

The evaluation unit is one validated source--target query over one generated Java class. Each retained query has an oracle Boolean label, an oracle canonical trace, and metadata describing the generation factors. The frozen release will contain \texttt{[N\_CLASSES]} generated classes and \texttt{[N\_INSTANCES]} queries. It will also record the final context-size grid, trace-length range, question count per retained trace length, guard settings, parameter settings, and random seeds as \texttt{[DATASET\_CONFIGURATION]}.

The dataset is produced by the benchmark generator rather than sampled from repositories. In the current generator, the default full sweep is restricted to Java and tree-shaped programs. The generator materializes each context as a directory containing the Java class, the prompt context, reachability questions, oracle chains, method names, and configuration metadata. For each requested trace length, it selects positive and negative queries in balanced pairs when both are available; depths that cannot supply enough valid pairs contribute fewer instances or are omitted from the retained release.

\begin{table*}[t]
\centering
\caption{Detailed configuration of the evaluation set. Both task families use the same generation grid.}
\label{tab:dataset-spec}
\scriptsize
\begin{tabularx}{\textwidth}{p{0.22\textwidth}p{0.24\textwidth}X}
\toprule
Category & Parameter & Value \\
\midrule
Dataset size & Java classes & 1,232 per task family \\
& Questions & 46,656 per task family \\
& Question polarity & 23,328 reachable and 23,328 unreachable \\
\midrule
Graph generation & Target class size & $\{100,250,500\}$ methods \\
& Absolute source--target distance & $\{1,4,7,10,13,16,19,22,25\}$ \\
& Local tree structure & Three trees of depth three; two calls per function \\
& Helper-call probability & 0.2 in total \\
\midrule
Generated code & Comments & $\{0,1\}$ \\
& Variables, loops, and parameters & $\{0,1,2\}$ \\
& Conditionals or guards & $\{0,1,3,5\}$ \\
\midrule
Factorial design & Code configurations & $3\times2\times3\times3\times4\times3=648$ \\
& Questions per distance and configuration & Four reachable and four unreachable \\
\midrule
Task semantics & Unguarded & All generated direct calls contribute to reachability \\
& Condition-aware & Only calls whose generated guards evaluate to true contribute to reachability \\
\bottomrule
\end{tabularx}
\end{table*}


\subsection{Experimental Factors}

The primary factors are task family, signed canonical trace length, generated class size, query polarity, and model configuration. For condition-aware instances, we additionally analyze guard-related metadata such as the number of generated \texttt{if} statements and whether predicted transitions use calls that are syntactically visible but not executable. Other code-surface settings, including comments, local variables, loops, method parameters, and declaration order, are retained as metadata and used only when they are part of the frozen dataset configuration.

The unguarded and condition-aware families are compared at the level of matched factor distributions. Unless the final release explicitly pairs the same abstract graph and query across families, we do not treat family comparisons as paired observations.

\subsection{Model Serving and Inference}

We evaluate \texttt{[N\_MODELS]} locally served LLM configurations. A configuration is not just a checkpoint name: it includes the serving backend, precision, tensor parallelism, context limit, decoding parameters, output budget, and whether the deployment exposes a reasoning-budget control. Reporting at the configuration level is necessary because the same checkpoint can behave differently under different serving and budget settings.

\begin{table*}[!t]
\centering
\caption{Model configurations used in the evaluation. Checkpoint and serving information will be added after the final configurations are frozen.}
\label{tab:models}
\scriptsize
\begin{tabularx}{\textwidth}{p{0.25\textwidth}p{0.22\textwidth}p{0.22\textwidth}X}
\toprule
Configuration & Inference mode & Budget metadata & Reproducibility metadata \\
\midrule
\texttt{[MODEL\_CONFIG\_1]} & \texttt{[MODE\_1]} & \texttt{[BUDGET\_CONTROL\_1]} & checkpoint, size, context, license and source, serving settings \\
\texttt{[MODEL\_CONFIG\_2]} & \texttt{[MODE\_2]} & \texttt{[BUDGET\_CONTROL\_2]} & checkpoint, size, context, license and source, serving settings \\
\texttt{[MODEL\_CONFIG\_3]} & \texttt{[MODE\_3]} & \texttt{[BUDGET\_CONTROL\_3]} & checkpoint, size, context, license and source, serving settings \\
\texttt{[MODEL\_CONFIG\_4]} & \texttt{[MODE\_4]} & \texttt{[BUDGET\_CONTROL\_4]} & checkpoint, size, context, license and source, serving settings \\
\texttt{[MODEL\_CONFIG\_5]} & \texttt{[MODE\_5]} & \texttt{[BUDGET\_CONTROL\_5]} & checkpoint, size, context, license and source, serving settings \\
\bottomrule
\end{tabularx}
\end{table*}


The intended serving stack uses vLLM on the Jean Zay supercomputer operated by IDRIS~\cite{cnrs2024jeanzay}. The inference runner starts a local OpenAI-compatible vLLM server, resolves model-context and GPU-memory settings, reads generated experiment directories, and writes one result file per directory. Final serving values, including backend version, GPU type, precision, tensor-parallel size, concurrency, context limit, and memory-utilization settings, are recorded as \texttt{[SERVING\_CONFIGURATION]}.

For reasoning-budget experiments, the runner supports an explicit sweep over requested reasoning budgets. When thinking mode is enabled, it first requests a bounded reasoning segment, then continues from that segment with a separate final-answer budget. When thinking mode is disabled or unavailable, the same task is run as a direct answer generation with the configured answer budget. The final paper will report budget effects only for configurations where the serving stack actually exposes and enforces the requested reasoning-budget control.

\subsection{Prompting}

Each prompt contains the generated Java class, the source method, the target method, and the required JSON output schema. It does not expose oracle labels, oracle traces, generator metadata, or evaluator internals. The required object has two fields: \texttt{final\_answer}, a Boolean reachability prediction, and \texttt{trace}, an ordered list of method names.

The prompt template differs only where the task family requires it. The unguarded template asks the model to follow method calls from the queried source. The condition-aware template additionally instructs the model to account for generated values and guards when deciding whether a visible call executes. The complete prompt templates and any illustrative examples are frozen before inference and included in the anonymized artifact as \texttt{[PROMPT\_ARTIFACT]}.

\subsection{Output Parsing}

Model outputs are parsed after inference by a deterministic parser. The parser extracts the final valid JSON object when extra text appears before or after the answer. A case is marked \texttt{NO\_RESPONSE} when no valid object can be decoded, when \texttt{final\_answer} is missing or cannot be interpreted as a Boolean, or when generation ends before an extractable answer appears.

The \texttt{trace} field is expected to be a JSON list of method names. Normalization is deliberately limited: it may handle unambiguous formatting variants, but it must not invent method names or repair semantically invalid traces. All primary accuracy metrics use strict scoring, so \texttt{NO\_RESPONSE} counts as incorrect and is also reported as a separate failure mode.

\subsection{Metrics}

\begin{table*}[t]
\centering
\caption{Evaluation metrics used for final answers, traces, and budget behavior.}
\label{tab:metrics}
\begin{tabularx}{\textwidth}{p{0.25\textwidth}X}
\toprule
Metric & Definition \\
\midrule
Final-answer accuracy & Predicted \true/\false matches the oracle label \\
Exact-chain accuracy & Predicted method chain exactly matches the expected oracle chain \\
Trace accuracy & Predicted trace satisfies the stricter trace representation used by the evaluator \\
Trace completeness & Fraction of the expected trace recovered in order \\
Emitted reasoning length & Tokenizer-counted visible reasoning text extracted from the model response; not a runtime consumed-token counter \\
\bottomrule
\end{tabularx}
\end{table*}


The metrics are designed to separate answer correctness from trace behavior. Final-answer accuracy measures the Boolean label only. Exact canonical trace accuracy requires the full predicted trace to match the oracle trace. Trace completeness gives ordered-recall credit for oracle methods recovered in canonical order, starting from the queried source. Graph-valid exploration checks whether predicted transitions are executable and whether the trace satisfies the family-specific stopping condition. Syntactic-call bias is evaluated on condition-aware instances by detecting predicted transitions that follow visible calls excluded from the executable graph by generated guards.

\subsection{Statistical Analysis}

The final analysis reports uncertainty for all primary proportions and bounded trace-completeness scores. The exact confidence-interval procedure is \texttt{[CI\_METHOD]}, and the trace-completeness interval procedure is \texttt{[TRACE\_CI\_METHOD]}. Factor-level analyses use the frozen metadata rather than post-hoc labels inferred from source text.

For the main factor analysis, candidate predictors are task family, signed canonical trace length, class size, query polarity, guard count, model configuration, and selected interactions justified by the research questions. If multiple queries come from the same generated class, the analysis accounts for within-class dependence using \texttt{[DEPENDENCE\_MODEL]}. Model-by-factor breakdowns that are not part of the preregistered analysis plan are treated as exploratory.

\subsection{Baselines and Sanity Checks}

The LLM configurations are the systems under study. Non-LLM baselines are used to calibrate interpretation, not to compete with the oracle.

\begin{itemize}[leftmargin=*]
    \item \textbf{Random-label baseline.} Predicts according to the empirical label distribution and emits a minimal trace, calibrating Boolean accuracy under label balance.
    \item \textbf{Syntactic-call baseline.} Treats every visible call as executable, estimating how far source-level call visibility alone can go.
    \item \textbf{Shortest-executable-path baseline.} Uses the executable graph to answer reachability and emit a shortest path for positive queries, separating semantic reachability from the stricter canonical trace convention.
    \item \textbf{Oracle-traversal sanity check.} Recomputes labels and traces from the validated executable graph and should match the stored oracle by construction.
\end{itemize}

The final set of reported baselines is \texttt{[MAIN\_BASELINES]}. Any implementation-only sanity checks are documented in the artifact rather than treated as experimental systems.

\subsection{Reproducibility}

The anonymized artifact will contain the frozen dataset, generated Java programs, oracle metadata, prompt templates, inference configurations, raw or filtered model outputs when shareable, parser and metric implementations, statistical scripts, and dashboard materials. The review artifact location is \texttt{[ARTIFACT\_URL]}. When model-license or storage constraints prevent sharing raw outputs, the artifact records the restriction and provides parsed outputs sufficient to reproduce the reported tables.
